\section{Discussion and conclusions}

Post-processing uses measurement information more effectively but does not modify the
quantum state. The classifier ablation gives a methodological result: strong predictive
metrics do not guarantee lower factorization cost. Decoding differs because syndromes
contain structured error information, and the hybrid works where the analytical decoder
misses noise correlations, although it must also be compared with other analytical
decoders. A learned model's value thus depends on its input and on what the baseline
leaves unused; the two applications involve different tasks and data, so their difference
cannot be attributed only to when machine learning intervenes.

For \(N = 15\), the AQFT gives a separate positive result: fewer gates accompany higher
success. The order \(r = 4\) yields phases exactly represented by two control-register
bits, so removed rotations carry no useful information and cost little precision. For
\(N = 21\), where \(r = 6\), phases are not exact and arithmetic is more complex, so the
benefit does not transfer automatically. Approximation must be chosen for each instance
and noise regime, then tested on data separate from selection.

The conclusions remain limited to simulation: post-processing uses \(N = 15\) with
uniform, stationary noise; QEC uses separate memory circuits; and AQFT tests use one
offline calibration and few instances and phases. The work demonstrates neither complete
fault-tolerant Shor nor the resources for cryptographic-scale factorization. It provides a
quantitative, reproducible characterization of when the studied components help, to be
extended through hardware validation, independent variation of \(N\) and \(r\), and
integration of protected logical operations, decoding, and arithmetic optimization. A
common rule emerges: complexity is justified only if it exploits information unused by
the baseline or reduces circuit cost by more than the loss introduced.
